\input{common_packages.tex}

\begin{document}

\title{\textbf{Simple Command Sequencer Generator} \\
\large\textit{Autocoder User Guide}}
\date{}
\maketitle

\section{Description}

The purpose of this generator is to provide a user-friendly way of defining named command sequences for the Simple Command Sequencer component. The generator takes a YAML model file as input which declares a suite of sequences -- each a named, ordered list of sub-commands (and optional sleeps) -- and autocodes four things from it: the Ada sequence step tables, bundled with the suite's frame-pool size into a \texttt{Config} constant consumed by the sequencer's \texttt{Init} routine; a set of first-class per-sequence wrapper commands so that each sequence appears on the sequencer instance as its own typed, operator-callable command; an operator-side command builder package for those wrapper commands; and the record type of the instance's periodic summary packet.

Note the example shown in this documentation is used in the unit test of this component so that the reader of this document can see it being used in context. Please refer to the unit test code for more details on how this generator can be used.

\section{Schema}

The following pykwalify schema is used to validate the input YAML model. Model files must be named in the form \textit{optional\_name.assembly\_name.command\_sequences.yaml}, where \textit{optional\_name} is the specific name of this set of sequences (only needed if more than one suite is defined in the assembly) and \textit{assembly\_name} is the assembly the sequences will be used in. Generally this file is created in the same directory or near to the assembly model file. The schema is commented to show each available YAML key and what it accomplishes; even without knowing the specifics of pykwalify schemas, you should be able to glean the available structure from the file below.

\yamlcodef{../schemas/command_sequences.yaml}

\section{Example Input}

The following is an example command sequences input YAML file. It defines eleven sequences exercising the main features: statically-armed commands and static sleeps, typed argument records whose fields are dispatched as dynamic arguments to sub-commands (including through non-byte-aligned packed fields and validatable records), a dynamic sleep whose duration comes from the sequence argument, no-wait direct dispatch, a sequence that calls another sequence as a sub-sequence step, and sequences whose command response is statically deferred until completion. This example adheres to the schema shown in the previous section, and is commented to give clarification.

Note the optional top-level \texttt{with:} key, which lists additional Ada packages to ``with'' in the generated package. A step's \texttt{arg:} expression is emitted verbatim into the generated Ada, so any package it references that is not pulled in automatically -- such as an enum package needed to qualify an enumeration literal -- must be listed here.

\yamlcodef{../../test/test_assembly/example_sequences.test_assembly.command_sequences.yaml}

\section{Example Output}

The example input above produces the following Ada outputs. The package spec declares the \texttt{To\_Arg} buffer-padding helper, the per-dynamic-step Resolver functions, a step array constant per sequence, the \texttt{Sequences\_Table} constant (plus a stable \texttt{Sequences} access pointer), and the suite's \texttt{Config} constant, which bundles the table access with the suite's \texttt{num\_concurrent\_sequences}. The package body implements the Resolver functions: each validates the sequence's per-call argument buffer through the input type's generated \texttt{Validation.Valid}, deserializes it through \texttt{Serialization.From\_Byte\_Array}, and re-serializes the addressed field as the step's argument buffer (a sub-command's argument, or a \texttt{Packed\_U32} millisecond count for a dynamic sleep). Resolvers exist only for dynamic steps -- a suite of purely static sequences declares none, and its body degenerates to \texttt{pragma No\_Body;}.

The \texttt{Config} constant should be passed into the Simple Command Sequencer's \texttt{Init} procedure during assembly initialization, e.g.\ in the assembly YAML:

\begin{verbatim}
- type: Simple_Command_Sequencer
  name: My_Sequencer
  init:
    - "Config => Test_Assembly_Command_Sequences_Example_Sequences.Config"
\end{verbatim}

\adacodef{../../test/test_assembly/build/src/test_assembly_command_sequences_example_sequences.ads}

\adacodef{../../test/test_assembly/build/src/test_assembly_command_sequences_example_sequences.adb}

\section{Per-Sequence First-Class Commands}

In addition to the sequence tables above, this same model drives the autocoded synthesis of one first-class command per declared sequence on the Simple Command Sequencer instance. Discovery is build-path-driven: the sequencer's commands model queries the model database for every \texttt{*.command\_sequences.yaml} in scope, so a naked component build (no assembly) sees no sequences and produces only the four static commands (\texttt{Run\_Sequence}, \texttt{Kill\_All\_Sequences}, \texttt{Kill\_Frame}, \texttt{Set\_Summary\_Packet\_Period}), while an assembly build whose path includes a sequences suite adds one command per sequence. Each per-sequence command (e.g.\ \texttt{Sequence\_B}) takes the sequence's own \texttt{arg\_type} verbatim as its command argument (or no argument for an argless sequence) -- a user-written, normally-registered type, so no argument records are generated or registered anywhere in the pipeline. At run time the component copies the received argument buffer unchanged into a generic \texttt{Run\_Sequence\_Arg.T} and dispatches through the existing \texttt{Run\_Sequence} backbone using the sequence's runtime slot index; the response behavior comes from the sequence's static configuration in the autocoded table (see Sequence Response Behavior below). A companion generated package, \texttt{<suite>\_Commands}, carries the operator-side builder surface (per-sequence command id getters and \texttt{Command.T} constructors) for unit tests and other on-board callers -- split from the sequences package because the sequencer component itself never uses the builders. No per-sequence Ada needs to be hand-written:

\adacodef{../../test/test_assembly/build/src/test_assembly_command_sequences_example_sequences_commands.ads}

\section{Sequences Calling Sequences}

Because every sequence is a first-class command on the sequencer instance, a sequence step may invoke another sequence simply by targeting the sequencer's own synthesized command, e.g.\ \texttt{command: My\_Sequencer.Sequence\_C}. No special syntax or generator support is involved -- the step resolves like any other \texttt{Component.Command} reference, provided the sequencer instance is declared in the assembly. At run time the sub-sequence command is emitted out the sequencer's \texttt{Command\_T\_Send} connector, routed through the command router, and arrives back at the sequencer's \texttt{Command\_T\_Recv\_Async}, where it claims its own frame and executes concurrently with (and independently of) the calling sequence. If the calling sequence waits for command completion, the sub-sequence's command response (per its static \texttt{response\_behavior} configuration, see below) is what wakes it. Nesting depth is naturally capped by \texttt{Num\_Concurrent\_Sequences}: a sub-sequence call when no frame is available is rejected with a \texttt{No\_Frame\_Available} event and a Failure response to the calling frame.

\section{Sequence Response Behavior}

When the sequencer replies to the command that started a sequence is controlled by the \texttt{Sequence\_Response\_Behavior} enumeration. For the synthesized per-sequence commands it is static configuration: the optional \texttt{response\_behavior:} field on each sequence in the YAML, baked into the autocoded sequence table. The generic \texttt{Run\_Sequence} command instead carries the behavior as an argument, so a per-call override is always available through it:

\begin{itemize}
  \item \texttt{Send\_After\_Sequence\_Start} (the conventional default) -- the command response is sent immediately once the sequence has been accepted and started on a frame. This is the right choice for fire-and-forget operation and for sub-sequence calls where the caller should proceed as soon as the sub-sequence is running.
  \item \texttt{Send\_After\_Sequence\_Completion} -- the command response is withheld until the sequence ends, and its status reflects the outcome: \texttt{Success} when the final step completes, or \texttt{Failure} when the sequence is aborted by a failed sub-command, times out waiting for a sub-command response, fails to resolve or validate a dynamic argument, encounters an out-of-range sleep or timeout, or is halted by \texttt{Kill\_All\_Sequences} or \texttt{Kill\_Frame}. The sequencer captures the originating command's source and id per frame, so concurrent deferred sequences each reply to their own operator.
\end{itemize}

A sequence step may forward a field of the parent sequence's argument to the sub-command it dispatches, using the \texttt{Arg.field.subfield} syntax in the step's \texttt{arg:}. Each such reference produces a typed Resolver in the generated package: at run time the sequencer dispatches that Resolver, which first validates the per-call argument buffer through the input type's generated \texttt{Validation.Valid} (a failure aborts the sequence with an \texttt{Invalid\_Dynamic\_Command\_Argument} event), then deserializes it through the input type's safe \texttt{Serialization.From\_Byte\_Array}, navigates to the addressed field, and re-serializes it as the sub-command's argument buffer. The resolver's signature accepts only a \texttt{Basic\_Types.Byte\_Array} -- there is no \texttt{'Address} reinterpretation -- so the dynamic-argument machinery is type-safe by construction; a resolver cannot overlay an arbitrary record onto raw memory, only deserialize a byte array through the input type's generated deserializer.

\section{Sleep Steps}

A sequence step may be a sleep instead of a command: \texttt{sleep\_ms:} pauses the executing frame without issuing any bus traffic. Sleep is not a bus command; it is interpreted intrinsically by the sequencer frame state machine. The duration takes two forms. A \emph{static} sleep (\texttt{sleep\_ms: 3000}) is a plain integer number of milliseconds fixed at build time; the model bounds it to $0 \dots 2^{31}-1$ so it always fits an \texttt{Ada.Real\_Time.Time\_Span} by construction, leaving wake-time overflow of \texttt{Sys\_Time} arithmetic as its only run-time failure. A \emph{dynamic} sleep (\texttt{sleep\_ms: Arg.Settle\_Time}) instead references the sequence's per-call argument with the same \texttt{Arg[.field]} syntax as dynamic command arguments; the referenced field must be a \texttt{Packed\_U32.T}, resolved at execution time through the same validating Resolver machinery (a validation failure aborts the sequence with an \texttt{Invalid\_Dynamic\_Sleep\_Argument} event) and range-checked before scheduling. A sleep that fails its range check or overflows the wake-time computation ends the sequence with a \texttt{Sequence\_Out\_Of\_Range\_Sleep} event and a \texttt{Failure} status to any pending deferred reply -- continuing would leave the frame parked on a stale wake time.

\section{Summary Packet Type}

The sequencer emits a periodic \texttt{Summary\_Packet} containing one \texttt{Sequence\_Frame\_Summary} record per sequence frame, in frame order. Its layout therefore depends on the suite's \texttt{num\_concurrent\_sequences}, so the generator also autocodes the packet's ground/documentation record type, \texttt{<suite>\_summary\_record.record.yaml}, alongside the suite's other outputs. The flight software never imports this type -- the component fills the packet buffer directly from its frame array -- but ground tools and assembly documentation need it. At assembly load time the sequencer's packets model resolves an instance's \texttt{Summary\_Packet} type to the summary record of the suite named by the instance's \texttt{Config} init parameter. Because the frame pool and the packet layout both come from the suite model's single \texttt{num\_concurrent\_sequences} value, they can never disagree; the model also rejects at build time any suite whose frame summaries would not fit in one packet buffer (the same bound the \texttt{Num\_Concurrent\_Sequences\_Type} subtype enforces at \texttt{Init}). The example input above generates the following record model:

\yamlcodef{../../test/test_assembly/build/yaml/test_assembly_command_sequences_example_sequences_summary_record.record.yaml}

\end{document}
